\section{Exercise 4}

Chương trình đọc một số nguyên, đưa giá trị sang \texttt{\$a0} làm tham số cho \texttt{factorial} và kiểm tra hai trường hợp trước khi gọi. Ngưỡng 12 là giới hạn của số nguyên 32 bit có dấu: $12! = 479\,001\,600$ vẫn nằm dưới $2^{31} - 1 = 2\,147\,483\,647$, còn $13! = 6\,227\,020\,800$ đã vượt:

\begin{codeblock}[language=MIPS]
input:
  li $v0, 5             # syscall 5: đọc số nguyên
  syscall

  move $a0, $v0         # $a0 = n: tham số cho factorial
  bltz $a0, error       # n < 0: không có giai thừa
  bgt $a0, 12, overflow # n > 12: kết quả vượt 32 bit
\end{codeblock}

Kết quả trả về trong \texttt{\$v0} được lưu sang \texttt{\$t0} rồi in ra:

\begin{codeblock}[language=MIPS]
  jal factorial
  move $t0, $v0 # $t0 = n!

  li $v0, 1     # syscall 1: in số nguyên
  move $a0, $t0 # in n!
  syscall

  j exit
\end{codeblock}

Giai thừa viết đệ quy theo $n! = n \times (n-1)!$, dừng ở $0! = 1$. Sau lời gọi \texttt{jal factorial}, \texttt{\$a0} đã bị tầng dưới sửa, nên giá trị $n$ của tầng hiện tại phải lưu vào stack cùng với \texttt{\$ra}. Khi tầng dưới trả về, \texttt{\$v0} đang giữ $(n-1)!$, nhân với $n$ vừa khôi phục để được kết quả:

\begin{codeblock}[language=MIPS]
factorial:
  beqz $a0, base_case # 0! = 1

  addi $sp, $sp, -4   # mở slot cho $ra
  sw $ra, 0($sp)      # lưu địa chỉ trả về
  addi $sp, $sp, -4   # mở slot cho n
  sw $a0, 0($sp)      # lưu n

  addi $a0, $a0, -1   # tham số mới: n - 1
  jal factorial       # $v0 = (n-1)!

  lw $t0, 0($sp)      # lấy lại n
  addi $sp, $sp, 4    # đóng slot của n
  lw $ra, 0($sp)      # khôi phục địa chỉ trả về
  addi $sp, $sp, 4    # đóng slot của $ra

  mul $v0, $v0, $t0   # n! = n * (n-1)!
  jr $ra

base_case:
  li $v0, 1           # 0! = 1
  jr $ra
\end{codeblock}

Hai nhãn lỗi chỉ in thông báo tương ứng rồi nhảy tới \texttt{exit}:

\begin{codeblock}[language=MIPS]
error:
  li $v0, 4    # syscall 4: in chuỗi
  la $a0, msg1 # "A factorial for your number can not be found."
  syscall
  j exit

overflow:
  li $v0, 4    # syscall 4: in chuỗi
  la $a0, msg2 # "Number is too large for 32-bit integer."
  syscall
  j exit
\end{codeblock}

Nhập 5, chương trình in:

\begin{codeblock}[language={}]
120
\end{codeblock}

Nhập $-2$, giá trị âm rẽ sang nhãn \texttt{error}:

\begin{codeblock}[language={}]
A factorial for your number can not be found.
\end{codeblock}

Nhập 13, giá trị vượt ngưỡng nên rẽ sang nhãn \texttt{overflow}:

\begin{codeblock}[language={}]
Number is too large for 32-bit integer.
\end{codeblock}
